\section{Ideal-preserving factor branching}
\label{sec:ideal-branching}

% What Algorithm Branching loses: the exponents, and with them the ideal
Algorithm~\ref{alg:Branching} discards the exponents $e_{q,f}$ and hence only
describes $V(P)$ and $\sqrt{\langle P\rangle}$. But ideal membership,
multiplicities, and embedded components depend on $I=\langle P\rangle$ itself.

% What this section contributes
In this section, we keep $I$. We show that the SAT status of the branches,
which Algorithm~\ref{alg:Branching} already computes, certifies members of
$I$. Replacing generators by these members gives a smaller system for the
same ideal (Algorithm~\ref{alg:IdealReduce}). We then decompose $I$ into
primary components,
\[
  I=\mathfrak q_1\cap\cdots\cap\mathfrak q_t.
\]
The minimal primes of Theorem~\ref{thm:prime-branching} give the isolated
components, so a general primary decomposition runs only when $I$ has embedded
components, and they sort its components into isolated and embedded ones
(Algorithm~\ref{alg:PrimaryBranching}).

\subsection{Ideal reduction}
\label{subsec:sat-reduction}

% Notation shared by the section
For a finite system $Q$ and $h\in C(Q)$, let
$a_h=\max\{e_{q,h}:q\in Q,\ h\in C_q\}$ be the largest exponent of $h$.

% The rule: a factor chosen by every SAT branch lies in the ideal
Suppose that a factor $h$ lies in every SAT minimal hitting set of $Q$. Then
$h$ vanishes on all of $V(Q)$, and the Nullstellensatz puts some power of $h$
in $\langle Q\rangle$. Lemma~\ref{lem:sat-branches} shows that the power
$h^{a_h}$ suffices, so the factorization bounds the exponent without a
Gr\"obner basis of $\langle Q\rangle$.

% Worked example: the rule finds a member that interreduction misses
For example, $Q=\{xy,\,xz,\,yz-1\}$ is interreduced for every monomial order.
Its factor sets $\{x,y\}$, $\{x,z\}$, and $\{yz-1\}$ have the minimal hitting
sets $\{x,yz-1\}$ and $\{y,z,yz-1\}$. The second is UNSAT, since $y=z=0$
contradicts $yz=1$. Hence $x$ lies in every SAT minimal hitting set, and
$a_x=1$. Indeed, the cofactors $y$ and $z$ of $x$ have no common zero on
$V(Q)$, and multiplying $1=z\cdot y-(yz-1)$ by $x$ gives
\[
  x=z\cdot xy-x\,(yz-1)\in\langle Q\rangle.
\]

% The reduction loop
Algorithm~\ref{alg:IdealReduce} adds these powers to the generators, drops
the generators they divide, and repeats on the new system
(Figure~\ref{fig:ideal_reduce_flowchart}). We say that $\mathcal B$
certifies $V(H)=\emptyset$ when Algorithm~\ref{alg:Branching} with
$\mathcal B$ certifies every terminal branch of $H$ as UNSAT. A minimal
hitting set whose status is unknown, for example after a time limit, counts
as SAT, which can only shrink $\Delta$. If no minimal hitting set is
SAT, the intersection is $C(Q)$, and Lemma~\ref{lem:sat-branches} still
applies. On the example, the first round replaces $xy$ and $xz$ by $x$, and
the second round finds no new member and returns $\{x,yz-1\}$.

\begin{algorithm}[htbp]
  \caption{Ideal reduction}
  \label{alg:IdealReduce}
  \begin{algorithmic}[1]
    \Require A finite polynomial system $P$ and a branch solver $\mathcal B$
    \Ensure A finite generating set $Q$ of $\langle P\rangle$
    \State $Q\gets\Call{Interreduce}{P}$
    \Loop
    \State $\mathcal H_{\mathrm{sat}}\gets\{H\in\mathcal H(Q):\mathcal B
    \text{ does not certify }V(H)=\emptyset\}$
    \State $\Delta\gets\{h^{a_h}:h\in\bigcap\mathcal
    H_{\mathrm{sat}}\}\setminus Q$
    \If{$\Delta=\emptyset$}
    \State \Return $Q$
    \EndIf
    \State $Q\gets\Call{Interreduce}{\Delta\cup\{q\in Q:\text{no
          element of $\Delta$
    divides $q$}\}}$
    \EndLoop
  \end{algorithmic}
\end{algorithm}

% A flow chart of ideal reduction
\begin{figure}[pos=htbp]
  \centering
  \resizebox{0.62\textwidth}{!}{\begin{tikzpicture}[
    font=\small,
    every node/.style={align=center},
    blue box/.style={draw=blue!70!cyan, line width=1pt, rounded corners=6pt,
    fill=blue!50!cyan!4, inner sep=7pt},
    loop box/.style={draw=orange!85!black, line width=1pt, rounded corners=6pt,
    fill=yellow!40!orange!12, inner sep=7pt},
    decision/.style={draw=green!55!black, line width=1pt, diamond, aspect=1.8,
    fill=green!60!teal!8, inner sep=1pt},
    green box/.style={draw=green!55!black, line width=1pt, rounded corners=6pt,
    fill=green!60!teal!4, inner sep=7pt},
    red box/.style={draw=red!70, line width=1pt, rounded corners=6pt,
    fill=red!4, inner sep=7pt},
    return box/.style={draw=violet, line width=1pt, rounded corners=6pt,
    fill=violet!6, inner sep=7pt},
    arr/.style={-{Stealth[length=2.6mm]}, line width=1pt},
  ]

  % ---- initialization ----
  \node[blue box] (init) {%
    \textcolor{blue}{\bfseries Initialize}\\
  $Q\gets\Call{Interreduce}{P}$};

  % ---- classify every minimal hitting set ----
  \node[loop box, below=8mm of init] (classify) {%
    \textcolor{orange!75!black}{\bfseries Classify}\\
    Solve the terminal branches of\\
  each $H\in\mathcal H(Q)$ with $\mathcal B$};

  % ---- two outcome branches ----
  \node[red box, anchor=west] (unsat) at ($(classify.south)+(8mm,-10mm)$) {%
    \textcolor{red}{\bfseries UNSAT}\\
  Drop $H$};

  \node[green box, anchor=east] (sat) at ($(classify.south)+(-8mm,-10mm)$) {%
    \textcolor{green!45!black}{\bfseries SAT or unknown}\\
  Keep $H$ in $\mathcal H_{\mathrm{sat}}$};

  % ---- bottom rail ----
  \coordinate (raill) at ($(sat.south)+(0,-4mm)$);
  \coordinate (railm) at (raill -| classify.south);
  \coordinate (railr) at (raill -| unsat.south);

  % ---- collect the certified members ----
  \node[blue box, below=7mm of railm] (collect) {%
    \textcolor{blue}{\bfseries Collect ideal members}\\
  $\Delta\gets\{h^{a_h}:h\in\bigcap\mathcal H_{\mathrm{sat}}\}\setminus Q$};

  % ---- decision, return, and replacement ----
  \node[decision, below=7mm of collect] (empty) {$\Delta=\emptyset$?};
  \node[return box, right=12mm of empty] (return)
  {\textcolor{violet}{\bfseries Return $Q$}};
  \node[blue box, below=7mm of empty] (replace) {%
    \textcolor{blue}{\bfseries Replace}\\
    Add $\Delta$, drop the generators that\\
  an element of $\Delta$ divides, interreduce};

  % ---- arrows: pipeline ----
  \draw[arr] (init) -- (classify);
  \coordinate (outcome split) at ($(classify.south)+(0,-6mm)$);
  \draw[line width=1pt] (classify.south) -- (outcome split);
  \draw[arr] (outcome split)
  .. controls +(0,-2.2mm) and +(-2.2mm,0) .. ++(4mm,-4mm)
  -- (unsat.west);
  \draw[arr] (outcome split)
  .. controls +(0,-2.2mm) and +(2.2mm,0) .. ++(-4mm,-4mm)
  -- (sat.east);
  \draw[line width=1pt] (sat.south) -- (raill) -- (railr) -- (unsat.south);
  \draw[arr] (railm) -- (collect);
  \draw[arr] (collect) -- (empty);
  \draw[arr] (empty) -- node[above] {yes} (return);
  \draw[arr] (empty) -- node[right] {no} (replace);

  % ---- feedback: the new system is classified again ----
  \coordinate (fbcol) at ($(sat.west)+(-8mm,0)$);
  \coordinate (feedback) at (fbcol |- classify.west);
  \draw[line width=1pt, rounded corners=5pt]
  (replace.west) -- (replace.west -| fbcol) -- (feedback);
  \draw[arr] (feedback) -- (classify.west);

\end{tikzpicture}
}
  \caption{Ideal reduction of Algorithm~\ref{alg:IdealReduce}. Every factor
    $h$ chosen by all SAT minimal hitting sets enters the system as $h^{a_h}$,
    which replaces the generators it divides. The loop stops when no new power
  appears, and $\langle Q\rangle=\langle P\rangle$ throughout.}
  \label{fig:ideal_reduce_flowchart}
\end{figure}

% What the reduced system is good for
The output generates $\langle P\rangle$, with the powers that the rounds find
in place of the products they divide.
When a single minimal hitting set $H$ is SAT and every exponent is $1$, one
round gives $\langle P\rangle=\langle H\rangle$, as in the example. The
ideal-level and the variety-level descriptions then coincide, and a
Gr\"obner basis of $\langle P\rangle$ costs no more than one of the branch.

% Optimization: only the intersection is needed, so most solver calls
% can be skipped
Each round needs only the intersection of the SAT minimal hitting sets, so an
implementation need not classify every minimal hitting set. It keeps the
factors $F$ that lie in every SAT hitting set checked so far, starting from
$F=C(Q)$, skips a hitting set that contains all of $F$, and stops checking
once $F$ is empty. Each checked SAT hitting set replaces $F$ by its
intersection with $F$, and each skipped one contains $F$, so the final $F$ is
$\bigcap\mathcal H_{\mathrm{sat}}$. Two disjoint SAT hitting sets already
empty $F$, so a round that finds nothing often costs only a few solver calls.

% Nonzero conditions: the rounds keep the saturation instead of the ideal
A polynomial $u$ that must not vanish, such as $\kappa_h$ in
Subsection~\ref{subsec:physical-system}, changes only the SAT test: a minimal
hitting set $H$ counts as SAT when $V(H)\not\subseteq V(u)$. The proof of
Lemma~\ref{lem:sat-branches} then shows that $Q$ and the cofactors of $h$
have no common zero outside $V(u)$, so the Nullstellensatz gives
$u^sh^{a_h}\in\langle Q\rangle$ for some $s\geq0$. Each round then keeps
$\langle Q\rangle:u^\infty$ instead of $\langle Q\rangle$, and the termination
argument of Theorem~\ref{thm:ideal-reduce} is unchanged.
Algorithm~\ref{alg:PrimaryBranching} then decomposes
$\langle Q\rangle:u^\infty$, whose minimal primes are those of
$\langle Q\rangle$ that do not contain $u$.

\subsection{Primary factor branching}
\label{subsec:primary-decomposition}

% What the earlier variants leave open
Algorithm~\ref{alg:Branching} with $\mathcal B_{\min}$ describes only
$\sqrt I$. A minimal primary decomposition of $I$
(Subsection~\ref{subsec:decompositions}) also records the multiplicity of
each component and the embedded components.

% What this variant contributes
Algorithm~\ref{alg:PrimaryBranching}
(Figure~\ref{fig:primary_branching_flowchart}) computes the minimal
primary decomposition. Let
$\mathcal B_{\mathrm{prim}}$ be an algorithm that returns a minimal primary
decomposition of its input, for
instance~\cite{gianni1988primary,shimoyama1996localization}.

% Step 1: a smaller system for the same ideal
% What each step contributes
The algorithm first replaces $P$ by the output $Q$ of
Algorithm~\ref{alg:IdealReduce}. By Theorem~\ref{thm:ideal-reduce},
$\langle Q\rangle=I$, so every later step works on $I$, and
$\mathcal B_{\mathrm{prim}}$ starts from a smaller input when
Algorithm~\ref{alg:IdealReduce} replaces products by powers of their factors.

% Step 2: the minimal primes
% The reduced system gives a second route to the same primes
% Optimizations: factor branching for the minimal primes
It then computes the minimal primes $\mathfrak p_1,\ldots,\mathfrak p_r$ of
$I$ with Algorithm~\ref{alg:Branching} on $Q$, the branch solver
$\mathcal B_{\min}$, and the postprocessing of
Subsection~\ref{subsec:minimal-prime-decomposition}. The minimal hitting sets
of $Q$ that the last round of Algorithm~\ref{alg:IdealReduce} certified as
UNSAT need not be revisited.

% Optimization: isolated components by separators, without the general solver
Optionally, the minimal primes give the isolated components without
$\mathcal B_{\mathrm{prim}}$. For each $\mathfrak p_i$, a \emph{separator} is
a polynomial $w_i\in\bigcap_{j\neq i}\mathfrak p_j\setminus\mathfrak p_i$.
One exists because $\mathfrak p_i$ contains no other minimal prime and hence,
by Subsection~\ref{subsec:decompositions}, not their intersection. Saturating
by $w_i$ removes every component whose radical does not contain
$\mathfrak p_i$~\cite{shimoyama1996localization}, and the top-dimensional
part of $I:w_i^\infty$, the intersection of its primary components of maximal
dimension~\cite{eisenbud1992direct}, is the isolated
$\mathfrak p_i$-primary component $\mathfrak q_i$
(Lemma~\ref{lem:isolated-separator}). Both computations start from a
Gr\"obner basis $G$ of $I$. If $\mathfrak q_1\cap\cdots\cap\mathfrak q_r=I$,
which $G$ decides, then $I$ has no embedded components, and the algorithm
returns $\mathfrak q_1,\ldots,\mathfrak q_r$ without running
$\mathcal B_{\mathrm{prim}}$.

% Step 3: one general primary decomposition of the reduced system
Otherwise, $\mathcal B_{\mathrm{prim}}$ decomposes $\langle Q\rangle$, which
is $I$.

% Step 4: the minimal primes sort the components
A component $\mathfrak q$ is isolated if $\sqrt{\mathfrak q}$ is one of the
$\mathfrak p_i$ and embedded otherwise. The radical of an embedded component
strictly contains at least one $\mathfrak p_i$, so the variety of the component
lies inside an irreducible component of $V(I)$.

% Worked example: an embedded component
For example, $I=\langle x^2,xy\rangle$ has the single minimal prime
$\langle x\rangle$, which needs no separator. The top-dimensional part of $I$
is $\langle x\rangle\neq I$, so $\mathcal B_{\mathrm{prim}}$ runs and returns
$I=\langle x\rangle\cap\langle x^2,y\rangle$. The component $\langle x\rangle$
is isolated, and $\langle x^2,y\rangle$ is embedded, since its radical
$\langle x,y\rangle$ strictly contains $\langle x\rangle$.

\begin{algorithm}[htbp]
  \caption{Primary factor branching}
  \label{alg:PrimaryBranching}
  \begin{algorithmic}[1]
    \Require A finite polynomial system $P$, branch solvers $\mathcal B$ and
    $\mathcal B_{\min}$, and a minimal primary decomposition algorithm
    $\mathcal B_{\mathrm{prim}}$
    \Ensure A minimal primary decomposition of $\langle P\rangle$, each
    component marked isolated or embedded
    \State $Q\gets\Call{IdealReduce}{P,\mathcal B}$
    \Comment{Algorithm~\ref{alg:IdealReduce}, $\langle
    Q\rangle=\langle P\rangle$}
    \State $\{\mathfrak p_1,\ldots,\mathfrak p_r\}\gets$ the minimal primes of
    $\langle Q\rangle$ \Comment{Theorem~\ref{thm:prime-branching}}
    \ForAll{$i\in\{1,\ldots,r\}$} \Comment{Optional}
    \State $w_i\gets$ a separator of $\mathfrak p_i$
    \State $\mathfrak q_i\gets$ the top-dimensional part of
    $\langle Q\rangle:w_i^\infty$
    \EndFor
    \If{$\mathfrak q_1\cap\cdots\cap\mathfrak q_r=\langle Q\rangle$}
    \Comment{Optional}
    \State \Return $\{\mathfrak q_1,\ldots,\mathfrak q_r\}$, all marked
    isolated
    \EndIf
    \State $\mathcal D\gets\mathcal B_{\mathrm{prim}}(Q)$
    \ForAll{$\mathfrak q\in\mathcal D$}
    \State mark $\mathfrak q$ isolated if
    $\sqrt{\mathfrak q}\in\{\mathfrak p_1,\ldots,\mathfrak p_r\}$,
    and embedded otherwise
    \EndFor
    \State \Return $\mathcal D$
  \end{algorithmic}
\end{algorithm}

% A flow chart of primary factor branching
\begin{figure}[pos=htbp]
  \centering
  \resizebox{0.79\textwidth}{!}{\begin{tikzpicture}[
    font=\small,
    every node/.style={align=center},
    blue box/.style={draw=blue!70!cyan, line width=1pt, rounded corners=6pt,
    fill=blue!50!cyan!4, inner sep=7pt},
    loop box/.style={draw=orange!85!black, line width=1pt, rounded corners=6pt,
    fill=yellow!40!orange!12, inner sep=7pt},
    decision/.style={draw=green!55!black, line width=1pt, diamond, aspect=1.8,
    fill=green!60!teal!8, inner sep=1pt},
    return box/.style={draw=violet, line width=1pt, rounded corners=6pt,
    fill=violet!6, inner sep=7pt},
    arr/.style={-{Stealth[length=2.6mm]}, line width=1pt},
  ]

  % ---- ideal reduction and the minimal primes ----
  \node[blue box, text width=45mm] (reduce) {%
    \textcolor{blue}{\bfseries Ideal reduction}\\
    $Q\gets\Call{IdealReduce}{P,\mathcal B}$\\
  $\langle Q\rangle=\langle P\rangle$};

  \node[blue box, text width=45mm, right=10mm of reduce] (primes) {%
    \textcolor{blue}{\bfseries Minimal primes}\\
    $\mathfrak p_1,\ldots,\mathfrak p_r$ of $\langle Q\rangle$ by\\
  Algorithm~\ref*{alg:Branching} with $\mathcal B_{\min}$};

  % ---- isolated components from separators (optional) ----
  \node[loop box, text width=45mm, below=8mm of primes] (separate) {%
    \textcolor{orange!75!black}{\bfseries Isolated components (optional)}\\
    $\mathfrak q_i\gets$ top-dimensional part of\\
  $\langle Q\rangle:w_i^\infty$, $w_i$ a separator of $\mathfrak p_i$};

  \node[decision, below=7mm of separate] (unmixed)
  {$\bigcap_i\mathfrak q_i=\langle Q\rangle$?};
  \node[return box, text width=45mm] (isolated) at (reduce |- unmixed) {%
    \textcolor{violet}{\bfseries Return $\mathfrak q_1,\ldots,\mathfrak q_r$}\\
  all isolated};

  % ---- the general decomposition, marked by the minimal primes ----
  \node[blue box, text width=45mm, below=7mm of unmixed] (decompose) {%
    \textcolor{blue}{\bfseries Primary decomposition}\\
    $\mathcal D\gets\mathcal B_{\mathrm{prim}}(Q)$; $\mathfrak
    q\in\mathcal D$ is\\
    isolated if $\sqrt{\mathfrak q}\in\{\mathfrak p_1,\ldots,\mathfrak p_r\}$,\\
  embedded otherwise};

  \node[return box, text width=45mm] (done) at (reduce |- decompose) {%
    \textcolor{violet}{\bfseries Return $\mathcal D$}\\
  minimal primary decomposition\\ of $\langle P\rangle$};

  % ---- arrows ----
  \draw[arr] (reduce) -- (primes);
  \draw[arr] (primes) -- (separate);
  \draw[arr] (separate) -- (unmixed);
  \draw[arr] (unmixed) -- node[above] {yes} (isolated);
  \draw[arr] (unmixed) -- node[right] {no} (decompose);
  \draw[arr] (decompose) -- (done);

\end{tikzpicture}
}
  \caption{Primary factor branching of Algorithm~\ref{alg:PrimaryBranching}.
    Ideal reduction shrinks the system to $Q$ without changing the ideal. The
    separators of the minimal primes give the isolated components, which are
    returned if they intersect to $\langle Q\rangle$. Otherwise the general
    solver $\mathcal B_{\mathrm{prim}}$ decomposes $Q$, and the minimal primes
  mark each component isolated or embedded.}
  \label{fig:primary_branching_flowchart}
\end{figure}

% Optimization: a separator from one saturation
A separator also needs no intersection of primes. The saturation
$I:\mathfrak p_i^\infty=\{g\in R:\mathfrak p_i^mg\subseteq I\text{
for some }m\geq0\}$
keeps exactly the components whose radicals do not
contain $\mathfrak p_i$, so it lies in every other minimal prime but not in
$\mathfrak p_i$. Hence every element of its Gr\"obner basis that does not
reduce to zero modulo a Gr\"obner basis of $\mathfrak p_i$ is a separator.
Which of the two constructions is cheaper depends on the ideal.

% Optimization: zero-dimensional ideals need no separators and have
% no embedded components
A zero-dimensional $I$ needs no separators. Its components are pairwise
comaximal, so $\dim_K R/I$ is the sum of the $\dim_K R/\mathfrak q_i$, each at
least $\dim_K R/\mathfrak p_i$ with equality only if
$\mathfrak q_i=\mathfrak p_i$. If
$\dim_KR/I=\sum_i\dim_KR/\mathfrak p_i$, every component is therefore prime.
Otherwise $\mathfrak q_i=I+\mathfrak p_i^m$ for the first $m$ at which
$\dim_K R/(I+\mathfrak p_i^m)$ stops growing, by Nakayama's
lemma~\cite[Prop.~2.6]{atiyah1969introduction}. All primes of a
zero-dimensional ideal are maximal, so it has no embedded components, and
the intersection test can be skipped.

% Implementation: comparing a radical with the minimal primes
The mark needs no equality test of ideals. A prime that strictly contains
$\mathfrak p_i$ has smaller dimension
(Subsection~\ref{subsec:minimal-prime-decomposition}), so $\mathfrak q$ is
isolated exactly when $\sqrt{\mathfrak q}$ contains some $\mathfrak p_i$ of its
own dimension. Containment means that the generators of $\mathfrak p_i$ reduce
to zero modulo a Gr\"obner basis of $\sqrt{\mathfrak q}$.

% Implementation: checking the result
Two properties of the output can also be checked independently. The
components must intersect to $I$, which a Gr\"obner basis of
$\langle P\rangle$ decides, and the isolated components must have exactly
$\mathfrak p_1,\ldots,\mathfrak p_r$ as their radicals.

\subsection{Theoretical analysis}
\label{subsec:ideal-branching-analysis}

\begin{lemma}[SAT branches]
  \label{lem:sat-branches}
  Let $\mathcal H_{\mathrm{sat}}\subseteq\mathcal H(Q)$ contain every
  SAT minimal
  hitting set, and let $h\in C(Q)$ lie in every $H\in\mathcal
  H_{\mathrm{sat}}$. Then
  $h^{a_h}\in\langle Q\rangle$.
\end{lemma}

\begin{proof}
  For $q\in Q$ with $h\in C_q$, let $r_q=q/h^{e_{q,h}}$ be the cofactor of $h$.
  We first show that the cofactors have no common zero on $V(Q)$. Suppose that
  $b\in V(Q)$ is a common zero of the cofactors. Then every $q\in Q$ has a
  factor $f\neq h$ with $f(b)=0$: a factor of $r_q$ if $h\in C_q$, and any
  factor that vanishes at $b$ otherwise. So the factors $f\neq h$ in $C(Q)$
  that vanish at $b$ form a hitting set, and a minimal hitting set $H$ inside
  it has $b\in V(H)$ and $h\notin H$. Then $H$ is SAT, hence
  $H\in\mathcal H_{\mathrm{sat}}$ and $h\in H$, a contradiction. By the
  Nullstellensatz, $1=\sum_{q\in Q}g_q\,q+\sum_{q:\,h\in C_q}g'_q\,r_q$ for
  some $g_q,g'_q\in R$. Since $e_{q,h}\leq a_h$, every
  $h^{a_h}r_q=h^{a_h-e_{q,h}}q$ lies in $\langle Q\rangle$, and multiplying by
  $h^{a_h}$ gives $h^{a_h}\in\langle Q\rangle$.
\end{proof}

\begin{theorem}[Ideal reduction]
  \label{thm:ideal-reduce}
  Algorithm~\ref{alg:IdealReduce} terminates and returns a finite
  generating set $Q$ of $\langle P\rangle$.
\end{theorem}

\begin{proof}
  Every round keeps $\langle Q\rangle$. Indeed, $\mathcal B$ certifies
  $V(H)=\emptyset$ only when it holds (Lemma~\ref{lem:recursive-refinement}), so
  $\mathcal H_{\mathrm{sat}}$ contains every SAT minimal hitting set. Hence
  $\Delta\subseteq\langle Q\rangle$ by Lemma~\ref{lem:sat-branches}, the dropped
  generators are multiples of elements of $\Delta$, and interreduction keeps the
  ideal.

  For termination, let $M(G)=\langle\operatorname{LM}(g):g\in G\rangle$ as in
  the proof of Theorem~\ref{thm:refine-termination}. Consider a round with
  $\Delta\neq\emptyset$, and let $Q_0=\Delta\cup\{q\in Q:\text{no
      element of $\Delta$ divides
  $q$}\}$, so that the round ends with $\Call{Interreduce}{Q_0}$. At
  the start of
  every round, $Q$ is interreduced, and in particular monic. Every $q\in Q$ is
  kept or is a multiple of an element of $\Delta$, so $M(Q)\subseteq
  M(Q_0)$. Let
  $\delta=h^{a_h}\in\Delta$ and let $q\in Q$ have $e_{q,h}=a_h$, so
  that $\delta\mid q$. If
  $\operatorname{LM}(\delta)\in M(Q)$, some $q'\in Q$ has
  $\operatorname{LM}(q')\mid\operatorname{LM}(\delta)\mid\operatorname{LM}(q)$,
  so
  $q'=q$ since $Q$ is interreduced. Then $q/\delta$ has leading
  monomial $1$, and
  $\delta=q\in Q$ since both are monic, which contradicts
  $\delta\notin Q$. Hence
  $M(Q)\subsetneq M(Q_0)$, and since interreduction does not shrink the
  leading-monomial ideal, $M(Q)$ grows strictly in every round that does not
  return. Since $R$ is Noetherian, there are finitely many rounds, and each
  makes finitely many calls of $\mathcal B$.
\end{proof}

\begin{lemma}[Isolated components by separators]
  \label{lem:isolated-separator}
  Let $\mathfrak p$ be a minimal prime of an ideal $I$ and $w$ a separator of
  $\mathfrak p$. Then $\mathfrak p$ is the only minimal prime of $I:w^\infty$,
  and the $\mathfrak p$-primary component of $I$ is the intersection of the
  primary components of $I:w^\infty$ of maximal dimension.
\end{lemma}

\begin{proof}
  Let $I=\mathfrak q_1\cap\cdots\cap\mathfrak q_t$ be a minimal primary
  decomposition with $\mathfrak p_j=\sqrt{\mathfrak q_j}$, and let
  $T=\{j:w\notin\mathfrak p_j\}$. By Subsection~\ref{subsec:decompositions},
  $I:w^\infty=\bigcap_{j\in T}\mathfrak q_j$, and this decomposition is again
  minimal, since a $\mathfrak q_j$ that contained the intersection of the other
  kept components would contain the smaller intersection of all other
  components. For $j\in T$, the prime $\mathfrak p_j$ contains $I$, hence
  $\sqrt I$, the intersection of the minimal primes of $I$, and hence one of
  them~\cite[Prop.~1.11]{atiyah1969introduction}. This minimal prime is
  $\mathfrak p$, since every other one contains $w$.
  The $\mathfrak p$-primary component $\mathfrak q$ of $I$ is kept, since
  $w\notin\mathfrak p$, so $\mathfrak p$ contains $I:w^\infty$. Every prime
  that contains $I:w^\infty$ contains some kept
  $\mathfrak q_j$~\cite[Prop.~1.11]{atiyah1969introduction}, hence
  $\mathfrak p_j\supseteq\mathfrak p$, so $\mathfrak p$ is the only minimal
  prime of $I:w^\infty$.

  It follows that $\mathfrak q$ is the isolated $\mathfrak p$-primary component
  of $I:w^\infty$, which is the same in every minimal primary decomposition of
  $I:w^\infty$. Every other component of such a decomposition has a radical
  that contains $I:w^\infty$, hence $\mathfrak p$, and differs from
  $\mathfrak p$. A prime that strictly contains $\mathfrak p$ has smaller
  dimension (Subsection~\ref{subsec:minimal-prime-decomposition}), so
  $\mathfrak q$ is the only component of maximal dimension.
\end{proof}

\begin{theorem}[Primary factor branching]
  \label{thm:primary-branching}
  Algorithm~\ref{alg:PrimaryBranching}, with or without its optional lines,
  terminates and returns a minimal primary decomposition $\mathcal D$ of
  $\langle P\rangle$:
  \[
    \langle P\rangle=\bigcap_{\mathfrak q\in\mathcal D}\mathfrak q.
  \]
  It marks a component isolated exactly when the component is an isolated
  component of $\langle P\rangle$, and the radical of every embedded component
  strictly contains a minimal prime of $\langle P\rangle$. If every component is
  given by a Gr\"obner basis, then $g\in\langle P\rangle$ if and only if $g$
  reduces to zero modulo the basis of every component.
\end{theorem}

\begin{proof}
  By Theorem~\ref{thm:ideal-reduce}, Algorithm~\ref{alg:IdealReduce} terminates
  and returns $Q$ with $\langle Q\rangle=\langle P\rangle=I$, so by
  Theorem~\ref{thm:prime-branching} the $\mathfrak p_i$ are the minimal primes
  of $I$. The remaining lines make finitely many solver calls, so the algorithm
  terminates.

  Suppose that the optional lines return. By
  Lemma~\ref{lem:isolated-separator}, each $\mathfrak q_i$ is the isolated
  $\mathfrak p_i$-primary component of $I$, and
  $\mathfrak q_1\cap\cdots\cap\mathfrak q_r=I$. The radicals $\mathfrak p_i$ are
  distinct, and no $\mathfrak q_i$ is redundant. Otherwise $\mathfrak p_i$
  contains $\bigcap_{j\neq i}\mathfrak q_j$, hence some $\mathfrak q_j$ with
  $j\neq i$~\cite[Prop.~1.11]{atiyah1969introduction}, and hence
  $\mathfrak p_j$, which contradicts the minimality of $\mathfrak p_i$. So the
  $\mathfrak q_i$ form a minimal primary decomposition of $I$, and all of them
  are isolated.

  Otherwise, $\mathcal B_{\mathrm{prim}}$ decomposes $\langle Q\rangle=I$. A
  component is isolated exactly when its radical is a minimal prime of $I$, that
  is, one of the $\mathfrak p_i$, which is the mark. The radical of an embedded
  component is a prime that contains $I$, hence some $\mathfrak p_i$ as in the
  proof of Lemma~\ref{lem:isolated-separator}, and it differs from
  $\mathfrak p_i$.

  In both cases, $g\in\langle P\rangle$ if and only if $g$ lies in every
  component, and a Gr\"obner basis decides membership by reduction
  (Subsection~\ref{subsec:groebner}).
\end{proof}
